\chapter{Explicit computations via NAP duality}
\label{chap:NAP-computations}

\index{Poincare duality!explicit computations}

\begin{remark}[Type signature; NAP as a level-$2$/$5$ statement]
\label{rem:nap-type-signature}
\index{type signature!NAP chapter}
The non-abelian Poincar\'e (NAP) chapter operates on the Open quadrant
in two strata. The bar/twisting calculation
$\barB^{\mathrm{ch}}_X(\cA)$ and its Verdier dual
$\cA^!_\infty$ are level-$2$ statements (bar coalgebra and post-Verdier
homotopy dual algebra). The genus-graded NAP scalar shadow
$\kappa(\cA)\lambda_g$ on $\overline{\cM}_g$ is a level-$5$ scalar
shadow, obtained by trace + clutching on the open factorization
category. The hypothesis package on every theorem in this chapter is:
\emph{finite type on the Koszul locus, or separated completed
continuous-dual replacement}; comparison maps from CY/Hall/Mukai data
\textup{(}Vol~III\textup{)} are not used as inputs and never identify
$\cA^!_\infty$ with a CY-side bulk. Modular invariance of any
genus-$g$ partition function specialization derived from NAP is a
level-$5$ trace+clutching statement; the closed shadow has modular
consequences but modularity is not an intrinsic closed-algebra
property. On the \emph{KSDual} sublocus
($\Z/2$-fixed under $\cA\mapsto\cA^!$) the Verdier-Koszul reflection
becomes an equivalence, and the genus complementarity of
Theorem~\ref{thm:genus-complementarity} reads as a Lagrangian
splitting rather than a comparison-with-comparison-map.
Cross-volume specializations enter at
$\Phi_d^{(\Sigma_{d-1}, C)}=\mathrm{Sp}^{\mathrm{ch}}_{\Sigma_{d-1},C}\circ
\Phi_d^{\mathrm{FA}}$ (Stage~$1$ at level~$1$, Stage~$2$ at level~$2$);
the K3-specific $C=E_\tau$ shadow appears in the Mukai-grading remark
of \S\ref{rem:nap-no-cy-hall-shortcut}.
\end{remark}

Non-abelian Poincar\'e duality enters this monograph only on the
Verdier-dualizable surface of Theorem~\ref{thm:main-NAP-resolution}:
finite type on the Koszul locus, or the separated completed
continuous-dual replacement. On this surface the chiral bar coalgebra
$\barB_X^{\mathrm{ch}}(\cA)$ is the factorization homology object,
and its Verdier dual
\[
\cA^!_\infty := \mathbb{D}_{\Ran}\barB_X^{\mathrm{ch}}(\cA)
\]
is the post-Verdier homotopy dual algebra. If the transferred
structure is formal and dualizable, the strict Koszul-dual algebra is
$\cA^!=(\cA^i)^\vee$, where
$\cA^i=H^\bullet\barB_X^{\mathrm{ch}}(\cA)$ is the Koszul-dual
coalgebra. The cobar counit
$\Omegach_X\barB_X^{\mathrm{ch}}(\cA)\simeq\cA$ is bar-cobar
inversion, not the definition of~$\cA^!$.

The computation is local on configuration strata. The logarithmic
forms on $\overline{C}_n(X)$ carry the collision residues; Verdier
duality turns the bar coproduct into the product of
$\cA^!_\infty$, with the usual dimension shift and orientation line.
Thus the same residue calculation detects the bar differential,
the post-Verdier curvature, and the scalar genus obstruction, but it
does so only after the finite-type or completed hypotheses have been
installed.

The examples isolate five mechanisms: the Heisenberg double pole,
where the bar-side curvature $+k$ becomes the post-Verdier curvature
$-k$; the free-fermion and $bc$-$\beta\gamma$ quadratic pairs, where
Koszul duality exchanges exterior and symmetric generators; the
affine family, where Verdier duality negates the shifted level
$k+h^\vee$; and the Virasoro family, where completion and the
Feigin--Frenkel involution produce the scalar shadow
$c\mapsto 26-c$.

\section{Worked examples: standard Koszul pairs}
\noindent
The NAP computations in this section are proved on the finite-type or
completed surfaces stated above. The critical-level affine-to-Yangian
bridge is isolated as Conjecture~\ref{conj:critical-affine-yangian}
and is not used as input.

\begin{remark}[The bar object is not a Hochschild complex]
\label{rem:nap-hochschild-separation}
The NAP complex in this chapter is the chiral bar coalgebra and its
Verdier dual. It is not
$Z^{\mathrm{der}}_{\mathrm{ch}}(\cA)
=\ChirHoch^\bullet(\cA,\cA)$, the curve-level chiral Hochschild
bulk of Theorem~H; it is not topological Hochschild homology of an
$E_n$-algebra; and it is not the categorical Hochschild complex of a
dg or Calabi--Yau category. Comparison maps may be constructed on
formal-disk or Stage--2 specialisation surfaces
(Chapter~\ref{chap:three-hochschild-unification}), but no equality
between these Hochschild theories is used here.
\end{remark}

\begin{remark}[No transitive proof through CY or Hall data]
\label{rem:nap-no-cy-hall-shortcut}
K3 Mukai pairings, Borcherds--Kac--Moody denominator weights, and
Hall--Drinfeld doubles live in their own categorical or Hall-algebra
ambients. They enter a curve-level NAP computation only through an
explicit comparison map
\[
D^b\mathrm{Coh}(Y)
\longrightarrow
\mathrm{Fact}^{\mathrm{hol}}(Y)
\longrightarrow
\mathrm{Fact}^{\mathrm{ch}}(X)
\longrightarrow
\barB_X^{\mathrm{ch}}(\cA)
\]
together with the induced Verdier pairing on configuration strata.
Absent that map, such data neither prove a Vol~I NAP claim nor
identify $\cA^i$, $\cA^!$, or
$Z^{\mathrm{der}}_{\mathrm{ch}}(\cA)$.
\end{remark}

\subsection{Heisenberg algebra}

\begin{example}[Heisenberg via NAP]\label{ex:heisenberg-NAP}
The Heisenberg algebra $\mathcal{H}_k$ with OPE
$J(z)J(w)=k/(z-w)^2$ has no simple-pole bracket, so
$d_{\mathrm{bracket}}=0$. The double pole gives the bar-side
curvature class $m_0^{\mathrm{bar}}=k\,\omega_{\mathrm{cen}}$
\textup{(}see~\S\ref{sec:frame-bar-all}\textup{)} and collision
residue $r^{\mathrm{coll}}_{\mathcal H}(z)=k\Omega_{\mathcal H}/z$
with rank-one coefficient $k/z$. The displayed strict dual is the
finite-rank current-space calculation; the oscillator Fock realization
uses the same conformal-weight completion and continuous dual
convention as the other free-field examples. Verdier duality reverses
the curvature sign, hence
\[
\mathcal{H}_k^!
\simeq
\bigl(\mathrm{Sym}^{\mathrm{ch}}(V^*),\,m_0^!=-k\,\omega_{\mathrm{cen}}\bigr)
\]
by Theorem~\ref{thm:heisenberg-not-self-dual}, and NAP gives
\[
\int_X \mathcal{H}_k
\simeq
\mathbb{D}\!\left(
\int_{-X}\bigl(\mathrm{Sym}^{\mathrm{ch}}(V^*),
m_0^!=-k\,\omega_{\mathrm{cen}}\bigr)\right).
\]
\end{example}

\begin{remark}[Heisenberg Koszul duality]\label{rem:heisenberg-koszul}
The duality $\mathcal{H}_k^! = \mathrm{Sym}^{\mathrm{ch}}(V^*)$ is \emph{not} level-rank duality $SU(N)_k \leftrightarrow SU(k)_N$, which exchanges different algebras. It is operadic: $\mathrm{chirLie}^! = \mathrm{chirCom}$, realized at the factorization algebra level by Verdier duality on configuration spaces.
\end{remark}

\subsection{Free fermions: correct Koszul pair}

\begin{example}[Free fermions via NAP]\label{ex:fermions-NAP}
The free fermion chiral algebra $\mathcal{F}$ has generators
$\psi(z), \psi^*(z)$ of conformal weight $h=1/2$ and OPE
\[\psi(z) \psi^*(w) = \frac{1}{z-w} + \text{regular}\]
\[\psi(z) \psi(w) = 0, \quad \psi^*(z) \psi^*(w) = 0\]

On the finite quadratic surface, or in the conformal-weight
completion of the Fock factorization algebra, the odd generator space
is
\[
V=\operatorname{span}\{\psi,\psi^*\}.
\]
As a graded algebra $\mathcal{F}\cong\Lambda(V)$, and classical
Koszul duality gives
\[
\Lambda(V)^! = \mathrm{Sym}(V^*).
\]

The chiral bar construction is represented on configuration strata by
\[
\bar{B}^{\mathrm{ch}}(\mathcal{F})
=
\bigoplus_{n \geq 0}
\int_{\overline{C}_{n+1}(X)}
\mathcal{F}^{\boxtimes(n+1)}\otimes\Omega^*_{\log}.
\]
The fermionic OPE $\psi \psi^* \sim (z-w)^{-1}$ gives the normalized
residue
\[
\operatorname{Res}_{D_{ij}}
(\psi_i \otimes \psi^*_j \otimes \eta_{ij})=1.
\]
The corresponding length-two coproduct is
\[\Delta(\psi^*) = 0, \quad \Delta(\psi) = 0 \quad \text{(primitives)}\]
\[\Delta(\psi \psi^*) = \psi \otimes \psi^* - \psi^* \otimes \psi\]
because the two generators are odd. After desuspension $s^{-1}$, the
cogenerators are even and the sign becomes the symmetric coalgebra
sign, giving $\mathrm{Sym}^c(V^*)$. Verdier duality identifies
\[
\int_X \mathcal{F}
\simeq
\mathbb{D}\!\left(\int_{-X}\beta\gamma\right),
\qquad
\mathcal{F}^! \simeq \beta\gamma .
\]
The algebra identification is Theorem~\ref{thm:fermion-boson-koszul}.
This is chiral Koszul duality
$\Lambda(V)^! \simeq \mathrm{Sym}(V^*)$ lifted to the vertex algebra
setting. The boson-fermion correspondence is a separate lattice-VOA
isomorphism $\mathcal{F} \cong V_{\mathbb{Z}}$, not this Koszul
duality.
\end{example}

\subsection{Affine Kac--Moody levels}

\begin{example}[Affine Lie algebras via NAP]\label{ex:kac-moody-NAP}
Let $\mathfrak{g}$ be simple with invariant form $(\, ,\,)$ and
dual Coxeter number $h^\vee$. The affine algebra
\[
\widehat{\mathfrak{g}}_k
=
(\mathfrak{g}\otimes\mathbb{C}[t,t^{-1}])\oplus\mathbb{C}K
\]
has OPE
\[
J^a(z)J^b(w)
\sim
\frac{k(a,b)}{(z-w)^2}
+\frac{[a,b](w)}{z-w}.
\]
On the generic completed finite-weight surface, the double pole gives
$r^{\mathrm{KM}}(z)=k\Omega/z$, while Verdier duality negates the
shifted curvature
\[
m_0(k)=\frac{k+h^\vee}{2h^\vee}\Omega .
\]
Thus the Koszul partner is
\[
\mathbb{D}_{\Ran}\barB^{\mathrm{ch}}(\widehat{\mathfrak{g}}_k)
\simeq
\widehat{\mathfrak{g}}_{k'},
\qquad
k'=-k-2h^\vee
\]
as in \eqref{eq:level-shifting-part1}; consequently
$\kappa(\widehat{\mathfrak g}_k)+
\kappa(\widehat{\mathfrak g}_{k'})=0$, while the ordinary central
charges satisfy
$c(\widehat{\mathfrak g}_k)+c(\widehat{\mathfrak g}_{k'})
=2\dim\mathfrak g$
\textup{(}Theorem~\ref{thm:central-charge-complementarity}\textup{)}.

At the critical level $k=-h^\vee$ the shifted curvature vanishes, the
level is fixed by $k\mapsto k'$, and the center becomes
infinite-dimensional:
\[
\mathcal{Z}(\widehat{\mathfrak{g}}_{-h^\vee})
\cong
\operatorname{Fun}(\operatorname{Op}_{{}^L\fg}(D^\times))
\]
by the Feigin--Frenkel theorem. The NAP statement available here is
the orientation-reversal statement on the completed center-aware
surface:
\[
\int_X \widehat{\mathfrak{g}}_{-h^\vee}
\simeq
\mathbb{D}\!\left(\int_{-X}\widehat{\mathfrak{g}}_{-h^\vee}\right)
\]
for the same chiral algebra. The Langlands dual appears in the
description of the center by opers; it is not produced by the NAP
duality functor.

\begin{conjecture}[Critical-level affine \texorpdfstring{$\to$}{to} Yangian bridge;
\ClaimStatusConjectured]\label{conj:critical-affine-yangian}
Following Costello's 4d Chern--Simons framework, one expects
\[
(\widehat{\mathfrak{g}}_{-h^\vee})^! \;\simeq\; Y(\mathfrak{g}).
\]
\end{conjecture}

\begin{remark}[Critical bridge is not an input]
Conjecture~\ref{conj:critical-affine-yangian} is not used as input in
this chapter. It records the missing comparison map from the completed
critical affine factorization algebra to the quantum-group interface.
\end{remark}
\end{example}

\subsection{Virasoro algebra}

\begin{example}[Virasoro via NAP]\label{ex:virasoro-NAP}
The Virasoro algebra has generator $T(z)$ with OPE:
\[T(z) T(w) = \frac{c/2}{(z-w)^4} + \frac{2T(w)}{(z-w)^2} + \frac{\partial T(w)}{z-w} + \text{reg}\]

The raw direct-sum bar complex is not the finite-type
pro-conilpotent surface: the OPE and harmonic corrections produce
infinite mode sums. The conformal-weight completed bar complex and
continuous Verdier dual are therefore part of the statement
\textup{(}Example~\ref{ex:virasoro-koszul-dual}\textup{)}.

The local collision kernel lowers OPE pole order by one. Hence
\[
r^{\mathrm{Vir}}_c(u)
=
\frac{c/2}{u^3}+\frac{2T}{u}
\]
\textup{(}Computation~\ref{comp:virasoro-rmatrix}\textup{)}, and the
scalar channel is the cubic-pole coefficient:
\[
\kappa(\mathrm{Vir}_c)=\frac{c}{2}.
\]
Equivalently, the Virasoro cocycle is
\[
\omega(L_m,L_n)=\frac{1}{12}m(m^2-1)\delta_{m+n,0}.
\]
The completed Koszul dual is curved, with scalar curvature
$m_0=(c/2)\omega_{\mathrm{Vir}}$.

At genus $g$ the completed bar construction is
\[
\bar{B}^{(g)}(\mathcal{V}_c)
=
\prod_{n \geq 0}
R\Gamma\!\left(
\overline{C}_n(\Sigma_g),
\mathcal{V}_c^{\boxtimes n}\otimes\Omega^*_{\log}
\right),
\]
with direct sum only in finite bar-length truncations.

The modular characteristic $\kappa(\mathrm{Vir}_c)=c/2$ couples to
the Hodge class on the uniform-weight scalar-diagonal lane of
Theorem~\ref{thm:genus-universality}:
\[
\mathrm{obs}_g(\mathrm{Vir}_c)
=
(c/2)\,\lambda_g
\in H^{2g}(\overline{\mathcal{M}}_g),
\]
and at genus~$1$ this reads $(c/2)\lambda_1$. Since Virasoro has one
strong generator, the cross-channel correction term of
Theorem~\ref{thm:genus-universality} vanishes for this family.

\begin{proposition}[Virasoro NAP duality at \texorpdfstring{$c=26$}{c=26};
\ClaimStatusConditional]\label{prop:virasoro-c26-selfdual}
\textbf{Type signature.} (Open quadrant; Verdier-Koszul reflection at
level~$2$; conformal-weight completed ambient required for
Theorem~\ref{thm:main-NAP-resolution}; $c=26$ is on the closed-shadow
modular consequence locus, not a self-dual KSDual locus.)

At $c = 26$, the Koszul dual is $\mathrm{Vir}_{26}^! \simeq \mathrm{Vir}_0$
(the uncurved algebra), so NAP gives
\[
\int_X \mathcal{V}_{26}
\;\simeq\;
\mathbb{D}\!\left(\int_{-X} \mathcal{V}_0\right).
\]
This is \emph{not} self-duality: the self-dual point under
the complementarity $c + c' = 26$
(Theorem~\textup{\ref{thm:central-charge-complementarity}})
is $c = 13$.
\end{proposition}

\begin{proof}
By Example~\ref{ex:virasoro-koszul-dual},
$\mathrm{Vir}_c^! \simeq \mathrm{Vir}_{26-c}$. Setting $c = 26$
gives $\mathrm{Vir}_{26}^! \simeq \mathrm{Vir}_0$.
The completed Verdier-dualizable form of
Theorem~\ref{thm:main-NAP-resolution} then gives the stated duality.
\end{proof}

\begin{remark}[Self-dual central charge as a KSDual locus]
\label{rem:c13-ksdual-poincare}
\index{KSDual sublocus!Virasoro $c=13$}
The self-dual point $c = 13$ is the \emph{KSDual} ($\Z/2$-fixed) locus
for the Virasoro family under $\cA \mapsto \cA^!$: it is the unique
central charge at which $\mathrm{Vir}_c^! \simeq \mathrm{Vir}_c$. On
this sublocus the Verdier-Koszul reflection at level~$2$ is an
equivalence rather than an adjunction with comparison map, and the
modular shadow on $\overline{\cM}_{1,n}$ \textup{(}level~$5$\textup{)}
carries the signature self-dual phase. Modular invariance of the
genus-$1$ trace on the open factorization category is recovered by
the trace + clutching package; on the KSDual sublocus the recovery
strengthens to an intrinsic $\Z/2$-equivariance, but it remains a
level-$5$ consequence of the open category, not an intrinsic
closed-algebra modularity claim. Whether the factorization homology
at $c = 13$ carries additional structure \textup{(}for example, a
non-degenerate self-pairing compatible with the completed Verdier
form\textup{)} is not proved here.
\end{remark}
\end{example}

\section{\texorpdfstring{Constructing $\mathcal{A}^!$ via non-abelian Poincar\'e duality}{Constructing the Koszul dual via NAP}}

\subsection{Residue construction}
\label{alg:NAP-koszul}

Given a chiral algebra $\mathcal{A}$ on $X$ satisfying the finite-type
Koszul hypotheses, or a separated completed replacement satisfying
continuous dualizability, let $\{a_i\}$ be homogeneous generators and
write the singular part of the OPE as
\[
a_i(z) a_j(w) = \sum_{k,m}\frac{C^k_{ij,m}}{(z-w)^m}a_k(w) + \text{descendants},
\]
the cohomological Koszul-dual chiral coalgebra $\mathcal{A}^i$
is constructed first; the Koszul-dual algebra is
$\mathcal{A}^!=(\mathcal{A}^i)^\vee$ when the finite-type or
completed continuous dual exists.

\emph{Cogenerators.}
For each generator $a_i \in \mathcal{A}$ introduce the desuspended
Verdier-linear dual cogenerator $a_i^*\in \mathcal{A}^i$. Its
cohomological degree is the dual degree shifted by the bar suspension;
its conformal weight is recorded separately and is fixed by the
residue pairing, not by the formula $|a_i^*|=-h_i$.

\emph{Coproduct from OPE residues.}
For each OPE term $a_i(z)a_j(w) \sim \frac{C^k_{ij,m}}{(z-w)^m} a_k(w)$, define the coproduct component:
\[
\left\langle \Delta(a_k^*),a_i\otimes a_j\right\rangle
=
(-1)^{\epsilon(i,j)}
\operatorname*{Res}_{z=w}
\left\langle a_k^*,a_i(z)a_j(w)\right\rangle ,
\]
where $\epsilon(i,j)$ is the Koszul sign obtained by moving the bar
desuspensions past the homogeneous fields. In coordinates this residue
extracts the coefficient $C^k_{ij,m}$ with the chosen logarithmic-form
normalisation.

\emph{Composite fields.}
If the OPE involves composite fields (e.g., $:T \cdot T:$ in W-algebras),
add generators $(\text{composite})^*$ to $\mathcal{A}^i$,
compute their coproducts from residue formulas,
and take the $I$-adic completion $\widehat{\mathcal{A}^i}$ before
dualizing.

\emph{Differential from the Verdier pairing.}
The coproduct
$\Delta\colon \mathcal{A}^i \to \mathcal{A}^i \otimes \mathcal{A}^i$ is
dual to the chiral product:
\[
\langle \Delta(a_i^*), a_j \otimes a_k \rangle = \langle a_i^*, \mu(a_j \otimes a_k)\rangle.
\]

\emph{Coalgebra axioms.}
Check coassociativity $(\Delta \otimes \mathrm{id}) \circ \Delta = (\mathrm{id} \otimes \Delta) \circ \Delta$,
the coderivation property $\Delta \circ d = (d \otimes \mathrm{id} + \mathrm{id} \otimes d) \circ \Delta$,
and conilpotency. In completed examples, such as Virasoro and
$W$-algebras, this last condition is topological conilpotency in the
weight or $I$-adic completion, not nilpotence in the raw direct-sum
coalgebra.

\emph{Koszul partner.}
If the Koszul dual coalgebra of $\mathcal{A}$ is quasi-isomorphic to the bar coalgebra of another chiral algebra $\mathcal{B}$, i.e.,
$\mathcal{A}^i\simeq \bar{B}^{\mathrm{ch}}(\mathcal{B})$,
then $(\mathcal{B}, \mathcal{A})$ form a Koszul pair. Verify:
\[
\Omega^{\mathrm{ch}}(\bar{B}^{\mathrm{ch}}(\mathcal{A})) \simeq \mathcal{A}
\quad\text{and}\quad
(\mathcal{A}^i)^\vee \simeq \mathcal{A}^!.
\]

\subsection{\texorpdfstring{Worked example: $\beta\gamma$ system}{Worked example: system}}

\begin{computation}[\texorpdfstring{$\beta\gamma$}{beta-gamma} Koszul dual]\label{comp:betagamma-dual}
The $\beta\gamma$ chiral algebra has generators $\beta(z), \gamma(z)$ with OPE:
\[\beta(z) \gamma(w) = \frac{1}{z-w} + \text{regular}\]
\[\beta(z) \beta(w) = 0, \quad \gamma(z) \gamma(w) = 0\]

On the finite quadratic surface, the underlying algebra is
$\mathrm{Sym}(V)$ for
$V=\operatorname{span}\{\beta,\gamma\}$; in the Fock model one works
with conformal-weight completions and continuous duals.

\emph{Dual generators.}
The coalgebra $(\beta\gamma)^i$ has desuspended dual cogenerators
$\beta^*,\gamma^*$. After Verdier-linear duality the strict algebra
$(\beta\gamma)^!$ is generated by the odd fields
\[
b\quad(h_b=\lambda),\qquad c\quad(h_c=1-\lambda).
\]

\emph{Coproduct from the simple pole.}

The OPE $\beta \gamma \sim (z-w)^{-1}$ gives:
\[\Delta(\beta^*) = 0 \quad \text{(primitive)}\]
\[\Delta(\gamma^*) = 0 \quad \text{(primitive)}\]

Since $\beta^*,\gamma^*$ are cogenerators of
$\bar{B}^{\mathrm{ch}}(\beta\gamma)$, the displayed primitive
coproduct is only the cogenerator part; the bar coalgebra also has
deconcatenation. The $\beta\gamma$ system is not a Koszul pair with
itself:
\[
(\beta\gamma)^!\simeq bc.
\]
This is quadratic chiral Koszul duality, not VOA bosonization.

\emph{Verdier verification.}

The Koszul-dual coalgebra is $\mathrm{co}\Lambda(V^*)$. Verdier
duality turns this coalgebra into the exterior algebra generated by
$b,c$. Thus
\[(\beta\gamma)^! = bc \quad \text{and} \quad (bc)^! = \beta\gamma\]

They are Koszul duals. The boson-fermion correspondence
$\mathrm{Sym}^*(V) \simeq \Lambda^*(V)$ as vector spaces is separate
from Koszul duality $\mathrm{Sym}^! = \Lambda^*$; the former is a
representation isomorphism, the latter an operadic duality. See
Theorem~\ref{thm:betagamma-fermion-koszul}.
\end{computation}

\section{Higher genus corrections via NAP}

\subsection{Genus expansion of factorization homology}

\begin{definition}[Genus-graded NAP complex]\label{def:genus-NAP}
\textbf{Type signature.} (Open quadrant; bar/twisting presentation on
$(\Sigma_g, D, \tau)$; level~$2$ at the bar coalgebra, level~$5$ for
the scalar $\kappa(\cA)\lambda_g$ projection; hypothesis package:
$\cA$ finite-type or completed Koszul on the relevant ambient.) The
modular shadow that emerges after trace + clutching on the open
factorization category is a level-$5$ consequence; it is not an
intrinsic closed-algebra invariant of $\cA$.

For a smooth genus-$g$ curve $\Sigma_g$ and a finite-type or completed
Koszul chiral algebra~$\mathcal{A}$, the genus-$g$ NAP chain model is
\[
\mathbf{C}^{\mathrm{NAP}}_g(\mathcal{A})
=
\prod_{n\geq0}
R\Gamma\!\left(
\overline{C}_n(\Sigma_g),
j_{n*}j_n^*(\mathcal{A}^{\boxtimes n})
\otimes\Omega^\bullet_{\log}
\right),
\]
with direct sum in the finite bar-length case. Its curved
differential is
\[
d_{\mathcal{A}}+d_{\mathrm{dR}}
+\sum_D \operatorname{Res}_D(\mu_{\mathcal{A}}),
\]
the sum over boundary collision divisors. Chiral associativity and
the Arnold relation reduce its square to the scalar curvature term.
After the flat representative is chosen, it becomes an ordinary
differential; on the uniform-weight scalar-diagonal lane
Theorem~\ref{thm:genus-universality} identifies the curvature with
$\kappa(\mathcal{A})\lambda_g$. In the all-weight setting the same
theorem supplies only the scalar summand
$\kappa(\mathcal{A})\lambda_g^{\mathrm{FP}}$ of the genus-$g$ free
energy; the cross-channel term
$\delta F_g^{\mathrm{cross}}(\mathcal{A})$ is additional data, not a
NAP consequence.
Verdier duality on the $n$-point stratum carries the usual
$[2n]$ orientation shift for a complex curve and sends this completed
bar chain model to the corresponding post-Verdier model for
$\mathcal{A}^!_\infty$.
\end{definition}

\begin{remark}[Curve Verdier shifts and Calabi--Yau Serre shifts]
\label{rem:curve-vs-cy-shifts}
The shift in Definition~\ref{def:genus-NAP} is the Verdier shift of
configuration spaces on a curve. It is not the Serre shift $[d]$ of a
$\mathrm{CY}_d$ input category. For a K3 input category
$D^b\mathrm{Coh}(K3)$ the categorical shift is $[2]$; an auxiliary
CY$_3$ factorization ambient used before Stage--2 specialisation does
not change the K3 Serre shift. After Stage--2 specialisation to a
reference curve, the resulting chiral algebra is governed by the
curve-level NAP complex above.
\end{remark}

\begin{remark}[Comparison maps from K3, BKM, and Hall--Drinfeld data]
\label{rem:k3-bkm-hall-comparison-map}
A Mukai lattice calculation, a BKM denominator identity, or a
Hall--Drinfeld double can supply a curve-level NAP input only after
the comparison map to $\mathbf{C}^{\mathrm{NAP}}_g(\mathcal{A})$ has
been constructed and checked against Verdier shifts, completions, and
residue normalisations. The K3 value in Theorem~C is a
Vol~III complementarity witness; it is not a proof of
Theorem~\ref{thm:main-NAP-resolution}, of
Definition~\ref{def:genus-NAP}, or of the examples in this chapter.
\end{remark}

\begin{theorem}[Genus complementarity; \ClaimStatusConditional]\label{thm:genus-complementarity}
\textbf{Type signature.} (Open quadrant; Verdier-Koszul reflection at
level~$2$, with the $\Z/2$-eigenspace splitting visible at level~$5$
on $\overline{\cM}_g$; hypothesis package: finite-type or completed
Koszul pair $(\cA,\cA^!)$, ordinary flat C0 representative, cochain
involution $\sigma$. The Lagrangian/dual splitting at $g\geq1$
specialises on the \emph{KSDual} sublocus
\textup{(}$\Z/2$-fixed points of $\cA\mapsto\cA^!$\textup{)} to the
fixed-point complex; off that sublocus, the splitting is the
adjunction-with-comparison-map decomposition.)

Let $(\mathcal{A},\mathcal{A}^!)$ be a finite-type or completed
Koszul pair for which the C0 ordinary flat representative of
Theorem~\ref{thm:quantum-complementarity-main} is available. Let
$\mathbf{C}_g(\mathcal{A})$ denote the ambient complementarity complex,
and assume the
Verdier--Koszul involution is represented there by a cochain
involution
\[
\sigma\colon \mathbf{C}_g(\mathcal{A})\to \mathbf{C}_g(\mathcal{A})
\]
as a map of complexes. Then
\[
\mathbf{C}_g(\mathcal{A})
\simeq
\mathbf{Q}_g(\mathcal{A})\oplus \mathbf{Q}_g(\mathcal{A}^!)
\]
with $\mathbf{Q}_g(\mathcal{A})$ and
$\mathbf{Q}_g(\mathcal{A}^!)$ the homotopy $\pm1$ eigenspaces of
$\sigma$. On cohomological shadows,
\[
H^*(\overline{\mathcal{M}}_g,\mathcal{Z}(\mathcal{A}))
=
Q_g(\mathcal{A})\oplus Q_g(\mathcal{A}^!).
\]
If $g\geq1$ and the Verdier pairing on
$\mathbf{C}_g(\mathcal{A})$ is non-degenerate and anti-invariant,
the two summands are Lagrangian and dual. At genus~$0$ the point
class is $\sigma$-fixed, so
$Q_0(\mathcal{A})=H^0(\overline{\mathcal{M}}_0,\mathcal{Z}(\mathcal{A}))$
and $Q_0(\mathcal{A}^!)=0$.
\end{theorem}

\begin{proof}
Theorem~\ref{thm:quantum-complementarity-main} gives the homotopy
eigenspace decomposition under exactly these C0 and
Verdier-involution hypotheses. The involution is constructed in
Lemma~\ref{lem:verdier-involution-moduli}; its compatibility with the
center action is the Kodaira--Spencer statement
Theorem~\ref{thm:kodaira-spencer-chiral-complete}. Since the base
field has characteristic zero, Lemma~\ref{lem:involution-splitting}
splits the complex into homotopy $\pm1$ eigenspaces. Taking
cohomology gives the displayed shadow decomposition on
$H^*(\overline{\mathcal{M}}_g,\mathcal{Z}(\mathcal{A}))$. The
Lagrangian and duality clauses are the positive-genus pairing clause
of Theorem~\ref{thm:quantum-complementarity-main}; the genus-$0$
exception is its final clause.
\end{proof}
